\chapter{計算実験}\label{computational_result}
実験結果を示す際には，計算環境（実験に用いた計算機のCPUやメモリー，
実装に用いた言語）を明記しましょう．
計算結果を表示するのに用いる図や表を表示する例を図\ref{fig1}と
表\ref{table1}に示しておきます．
また，今回は2段組ではありませんが, 2段組みの原稿において幅の広い表を表示する例を表\ref{table2}に示しておきます
（図も同様に\texttt{\char`\\begin\{figure*\}}のようにすればよい）．

\begin{figure}[htbp]
 \centering
 \includegraphics[width=0.9\linewidth]{sample_fig.eps}
 \caption{図の表示例}
 \label{fig1}
\end{figure}

\begin{table}[htbp]
 \centering
 \tabcolsep = 15pt
 \renewcommand{\arraystretch}{0.8}
 \caption{表の表示例}
 \label{table1}
 \begin{tabular}{lrr} \hline
  問題例 & 最良値 & 計算時間（秒） \\ \hline
  c05100 &    123 & 10.1 \\
  c10100 &    456 & 15.2 \\
  c20100 &    789 & 20.3 \\ \hline
 \end{tabular}
\end{table}


\begin{table*}
 \centering
 \tabcolsep = 19pt
 \renewcommand{\arraystretch}{0.8}
 \caption{2段組みスタイルにおいて幅の広い表を表示する例}
 \label{table2}
 \begin{tabular}{lrrcrr} \hline
  &\multicolumn{2}{c}{既存手法} && \multicolumn{2}{c}{提案手法}\\ \cline{2-3} \cline{5-6}
  
  問題例 & 最良値 & 計算時間（秒）&& 最良値 & 計算時間（秒） \\ \hline
  c05100 &    123 &          10.1 &&    111 &          10.0 \\
  c10100 &    456 &          15.2 &&    432 &          15.0 \\
  c20100 &    789 &          20.3 &&    765 &          20.0 \\ \hline
 \end{tabular}
\end{table*}